\documentclass[10pt,a4paper]{amsart}
\usepackage[margin=19mm]{geometry}
\usepackage{amsmath,amssymb,mathtools}
\usepackage[scr=rsfs,cal=euler]{mathalfa}
\usepackage{xfrac}
\usepackage[unicode,hidelinks,bookmarks=false]{hyperref}
\newcommand{\ie}{i.e.\ }
\newcommand{\cf}{cf.\ }
\newcommand{\resp}{respectively\ }
\newcommand{\Subsection}{\subsection}
\providecommand{\nobreakdash}{\nobreak\hbox{-}}
\setlength{\emergencystretch}{2em}
\multlinegap0pt
\usepackage{bm}
\usepackage{bbm}
\usepackage{dsfont}
\usepackage{xspace}
\usepackage{cancel}
\usepackage{mathtools}
\usepackage{amsthm}
\makeatletter
\@ifundefined{theorem}{\newtheorem{theorem}{Theorem}}{}
\@ifundefined{lemma}{\newtheorem{lemma}[theorem]{Lemma}}{}
\@ifundefined{proposition}{\newtheorem{proposition}[theorem]{Proposition}}{}
\makeatother
\makeatletter
\expandafter\ifx\csname customthm\endcsname\relax
\newenvironment{customthm}[1]
{\renewcommand\theinnercustomthm{#1}\innercustomthm}
{\endinnercustomthm}
\fi
\makeatother
\title[Co-rotating nearly parallel helical vortices with small cros]{Co-rotating nearly parallel helical vortices with small cross-section in 3D incompressible Euler equations}
\author{Daomin Cao, Jie Wan}
\date{Source: arXiv:2511.05956v1 (CC BY 4.0)}
\begin{document}
\maketitle

\noindent\emph{Source and licence.} Co-rotating nearly parallel helical vortices with small cross-section in 3D incompressible Euler equations. Version arXiv:2511.05956v1 (CC BY 4.0), distributed under the Creative Commons Attribution 4.0 International licence, as recorded in the arXiv licence metadata for this exact version and shown on the abstract page https://arxiv.org/abs/2511.05956v1. Full rights notice: licenses/arxiv-2511.05956-CC-BY-4.0.txt.\par\smallskip

\noindent\textit{Editorial scope.} Counted unit: the Proposition whose source label is \texttt{coercive esti} in the article (source0.tex lines 1166--1171, source label \texttt{coercive esti}), delivered together with the article's own complete original proof of it (source0.tex lines 1173--1396, one \texttt{proof} environment of 224 lines). No mathematics is written, repaired or re-derived: statement and proof are the article's own text line for line. Required context retained uncounted: 3-001 (lines 889--896); 203 (lines 989--993); 200 (lines 1020--1027); lemA-5 (lines 1032--1041); limit eq (lines 1107--1109); Non-degenerate (lines 1111--1113); coef of C (lines 1146--1152). Every definition, display and label of those lines is retained unchanged. Omitted from this package and not redistributed: 1--888, 897--988, 994--1019, 1028--1031, 1042--1106, 1110--1110, 1114--1145, 1153--1165, 1172--1172, 1397--3366. Nothing inside a retained excerpt is omitted or altered: every retained body is delivered exactly as the article's lines are written, so no omission receipt of any kind accompanies this package. Citations: the retained excerpts carry no citation command at all, so no bibliography entry of the article is transcribed and no reference is redistributed. Reference adaptation: none was needed and none was made; every reference inside the delivered unit resolves inside it, so no unresolved reference or citation is produced. Printed slips kept literal: no printed word, symbol, hypothesis or number of the counted statement or of its proof was altered. Layout and class: the article's own class is \texttt{amsart} with packages including amsmath, amssymb, amstext, a4wide, graphicx, color, cases, hyperref; the wrapper is the lane's pinned \texttt{amsart} editorial wrapper at 10pt on A4, which restates the theorem-like environments the retained text needs and adds the article's own macro definitions that the retained text needs (none), with extra wrapper packages (bm, bbm, dsfont, xspace, cancel, mathtools). The delivered numbering is the unit's own. Assets: no retained span contains a figure environment, a graphics-inclusion command, a PDF-page inclusion, a table or a TikZ picture; no image file is copied into this package, the package declares no asset dependency and no image file is redistributed with it. Licence: version arXiv:2511.05956v1 of this article is distributed under the Creative Commons Attribution 4.0 International licence, as recorded in the arXiv licence metadata for this exact version and shown on the abstract page https://arxiv.org/abs/2511.05956v1. A full rights notice is delivered at licenses/arxiv-2511.05956-CC-BY-4.0.txt. Characters: every retained excerpt is pure ASCII, so the delivered engine, XeLaTeX with its default Latin Modern text font, reports no missing character.
\begin{equation}\label{3-001}
||H_{\varepsilon, \hat{x}, \hat{q}}||_{W^{2,p}(\Omega)}\leq
\begin{cases}
\frac{C}{\varepsilon^{1-\frac{2}{p}}},\ \ &p>2,\\
C|\ln\varepsilon|,\ \ &p=2,\\
C,\ \ &1\leq p<2.
\end{cases}
\end{equation}

\begin{equation}\label{203}
\begin{split}
\sum_{j=1}^NV_{\varepsilon, Z,j}(x)-q(x)|\ln\varepsilon|=V_{\varepsilon, z_i, \hat{q}_{\varepsilon,i}}(x)-\hat{q}_{\varepsilon,i}|\ln\varepsilon|+O\left( \varepsilon^\gamma\right).
\end{split}
\end{equation}

\begin{equation}\label{200}
\begin{split}
\frac{\partial V_{\varepsilon, z_i, \hat{q}_{\varepsilon,i}}(x)}{\partial x_\hbar}=\begin{cases}
\frac{1}{s_{\varepsilon,i}}(\frac{\varepsilon}{s_{\varepsilon,i}})^{\frac{2}{p-1}}\phi'(\frac{|T_{z_i}(x-z_i)|}{s_{\varepsilon,i}})\frac{(T_{z_i})_\hbar^T\cdot T_{z_i}(x-z_i)}{|T_{z_i}(x-z_i)|},\ \ &|T_{z_i}(x-z_i)|\leq s_{\varepsilon,i},\\
\frac{\hat{q}_{\varepsilon,i}|\ln\varepsilon|}{\ln s_{\varepsilon,i}}\frac{(T_{z_i})_\hbar^T\cdot T_{z_i}(x-z_i)}{|T_{z_i}(x-z_i)|^2},\ \ &|T_{z_i}(x-z_i)|> s_{\varepsilon,i},
\end{cases}
\end{split}
\end{equation}

\begin{lemma}\label{lemA-5}
	Let $\gamma\in(0,1)$ and $ Z\in \Lambda_{\varepsilon, N} $. Then there exists a constant $ L > 1 $ such that for $ \varepsilon $ small
	\begin{equation*}
	V_{\varepsilon,Z}-q|\ln\varepsilon|>0,\ \ \  \ \text{in}\ \   \cup_{j=1}^N\left( T_{z_j}^{-1}B_{\left( 1-L \varepsilon^\gamma \right) s_{\varepsilon,j}}(0)+z_j\right),
	\end{equation*}
	\begin{equation*}
	V_{\varepsilon,Z}-q|\ln\varepsilon|<0,\ \ \  \ \text{in}\ \  \Omega\backslash\cup_{j=1}^N\left( T_{z_j}^{-1}B_{Ls_{\varepsilon,j}}(0)+z_j\right).
	\end{equation*}
	
\end{lemma}

\begin{equation}\label{limit eq}
-\Delta v-pw^{p-1}_+v=0, \ \ v\in L^{\infty}(\mathbb{R}^2).
\end{equation}

\begin{proposition}[Non-degeneracy]\label{Non-degenerate}
$ w $ is non-degenerate, i.e., the kernel of the linearized equation \eqref{limit eq} is $$ span\left \{\frac{\partial w}{\partial x_1}, \frac{\partial w}{\partial x_2}\right \}. $$
\end{proposition}

\begin{equation}\label{coef of C}
\begin{split}
\int_{\Omega}&\left( -\varepsilon^2\text{div}\left (K(z_j)\nabla   \frac{\partial V_{\varepsilon, z_j, \hat{q}_{\varepsilon,j}}}{\partial x_{\tilde{h}}}\right )\right)\left( \eta_i\frac{\partial V_{\varepsilon, Z,i}}{\partial x_\hbar}\right)\\
&=p\int_{\Omega}\eta_i\left (V_{\varepsilon, z_j, \hat{q}_{\varepsilon,j}}-\hat{q}_{\varepsilon,j}|\ln\varepsilon|\right )^{p-1}_+ \frac{\partial V_{\varepsilon, z_j, \hat{q}_{\varepsilon,j}}}{\partial x_{\tilde{h}}}\frac{\partial V_{\varepsilon, z_i, \hat{q}_{\varepsilon,i}}}{\partial x_\hbar}+O\left( \varepsilon^\gamma\right)\\
&=\delta_{i,j}(M_{i})_{\tilde{h},\hbar}+O\left( \varepsilon^\gamma\right) ,
\end{split}
\end{equation}

\begin{proposition}\label{coercive esti}
There exist  $ \rho_0>0, \varepsilon_0>0 $ such that for any $ \varepsilon\in(0,\varepsilon_0],  Z\in \Lambda_{\varepsilon, N} $, if $ u\in E_{\varepsilon,Z} $ satisfying $ Q_\varepsilon L_\varepsilon u=0 $ in $ \Omega\backslash\cup_{j=1}^NB_{Ls_{\varepsilon,j}}(z_j) $ for some $ L>1 $ large, then
\begin{equation*}
||Q_\varepsilon L_\varepsilon u||_{L^p}\geq \rho_0\varepsilon^{\frac{2}{p}}||u||_{L^\infty}. 
\end{equation*}
\end{proposition}

\begin{proof}
We argue by contradiction. Suppose that there are $ \varepsilon_m \to 0 $, $ Z_m=(z_{m,1},\cdots,z_{m,N})\to (z_{1},\cdots,z_{N})\in \Omega^{(N)} $  and $  u_m \in E_{\varepsilon_m, Z_m} $ with $ Q_{\varepsilon_m} L_{\varepsilon_m} u_m = 0 $ in $ \Omega \backslash\cup_{j=1}^NB_{Ls_{\varepsilon_m,j}}(z_{m,j}) $ for some $ L $ large and $ ||u_m||_{L^\infty}=1 $ such that
\begin{equation*}
||Q_{\varepsilon_m} L_{\varepsilon_m} u_m||_{L^p}\leq \frac{1}{m}\varepsilon_m^{\frac{2}{p}}.
\end{equation*}
By definition of $ Q_{\varepsilon}  $,
\begin{equation}\label{208}
Q_{\varepsilon_m} L_{\varepsilon_m} u_m=  L_{\varepsilon_m} u_m-\sum_{j=1}^N\sum_{\tilde{h}=1}^2b_{j,\tilde{h},m}\left( -\varepsilon_m^2\text{div}\left (K(z_{m,j})\nabla   \frac{\partial V_{\varepsilon_m, z_{m,j},\hat{q}_{\varepsilon_m,j}}}{\partial x_{\tilde{h}}}\right )\right) .
\end{equation}
We now estimate $ b_{j,\tilde{h},m} $. For fixed $ i=1,\cdots,N, \hbar=1,2 $, multiplying \eqref{208} by $ \eta_i\frac{\partial V_{\varepsilon_m, Z_{m},i}}{\partial x_\hbar} $ and integrating on $ \Omega $ we get
\begin{equation*}
\begin{split}
\int_{\Omega}u_m&L_{\varepsilon_m}\left( \eta_i\frac{\partial V_{\varepsilon_m, Z_{m},i}}{\partial x_\hbar}\right) =\int_{\Omega}L_{\varepsilon_m}u_m\left( \eta_i\frac{\partial V_{\varepsilon_m, Z_{m},i}}{\partial x_\hbar}\right) \\
=&\sum_{j=1}^N\sum_{\tilde{h}=1}^2b_{j,\tilde{h},m}\int_{\Omega}-\varepsilon_m^2\text{div}\left (K(z_{m,j})\nabla   \frac{\partial V_{\varepsilon_m, z_{m,j},\hat{q}_{\varepsilon_m,j}}}{\partial x_{\tilde{h}}}\right )\left( \eta_i\frac{\partial V_{\varepsilon_m, Z_{m},i}}{\partial x_\hbar}\right) .
\end{split}
\end{equation*}

%We estimate $ \int_{\Omega}u_NL_{\delta_N}\left( \eta_i\frac{\partial V_{\delta_N, Z_{N},i}}{\partial x_\hbar}\right). $ 
For the left hand side of the above equality, we have
\begin{equation*}
\begin{split}
&\int_{\Omega}u_mL_{\varepsilon_m}\left( \eta_i\frac{\partial V_{\varepsilon_m, Z_{m},i}}{\partial x_\hbar}\right) \\
=&-\int_{\Omega}u_m\varepsilon_m^2\text{div}\left (K(x)\nabla \left( \eta_i\frac{\partial V_{\varepsilon_m, Z_{m},i}}{\partial x_\hbar}\right) \right )- p\int_{\Omega}u_m\left (V_{\varepsilon_m, Z_m}-q|\ln\varepsilon_m|\right )^{p-1}_+\left( \eta_i\frac{\partial V_{\varepsilon_m, Z_{m},i}}{\partial x_\hbar}\right)  \\
=&-\int_{\Omega}\eta_iu_m\varepsilon_m^2\text{div}\left (K(x)\nabla \frac{\partial V_{\varepsilon_m, Z_{m},i}}{\partial x_\hbar} \right )-2\int_{\Omega}u_m\varepsilon_m^2\left (K(x)\nabla\eta_i|\nabla   \frac{\partial V_{\varepsilon_m, Z_{m},i}}{\partial x_\hbar}  \right )\\
&-\int_{\Omega}u_m\varepsilon_m^2\text{div}\left (K(x)\nabla   \eta_i  \right )\frac{\partial V_{\varepsilon_m, Z_{m},i}}{\partial x_\hbar}- p\int_{\Omega}u_m\left (V_{\varepsilon_m, Z_m}-q|\ln\varepsilon_m|\right )^{p-1}_+\left( \eta_i\frac{\partial V_{\varepsilon_m, Z_{m},i}}{\partial x_\hbar}\right)
\end{split}
\end{equation*}
\begin{equation}\label{301}
\begin{split}
=&\int_{\Omega}\eta_iu_mp\left( V_{\varepsilon_m, z_{m,i}, \hat{q}_{\varepsilon_m,i}}-\hat{q}_{\varepsilon_m,i}|\ln\varepsilon_m|\right)^{p-1}_+\frac{\partial V_{\varepsilon_m, z_{m,i}, \hat{q}_{\varepsilon_m,i}}}{\partial x_\hbar}+\int_{\Omega}\eta_iu_m\varepsilon_m^2\text{div}\left (\frac{\partial K(x)}{\partial x_\hbar}\nabla V_{\varepsilon_m, Z_{m},i}  \right )\\
&-2\int_{\Omega}u_m\varepsilon_m^2\left (K(x)\nabla\eta_i|\nabla   \frac{\partial V_{\varepsilon_m, Z_{m},i}}{\partial x_\hbar}  \right ) -\int_{\Omega}u_m\varepsilon_m^2\text{div}\left (K(x)\nabla   \eta_i  \right )\frac{\partial V_{\varepsilon_m, Z_{m},i}}{\partial x_\hbar}\\
&- p\int_{\Omega}u_m\left (V_{\varepsilon_m, Z_m}-q|\ln\varepsilon_m|\right )^{p-1}_+\left( \eta_i\frac{\partial V_{\varepsilon_m, Z_{m},i}}{\partial x_\hbar}\right).
%=:&I_1+I_2+I_3+I_4+I_5.
\end{split}
\end{equation}

By \eqref{3-001}, \eqref{203} and Lemma \ref{lemA-5},  
\begin{equation}\label{302}
\begin{split}
\int_{\Omega}\eta_i&u_mp\left( V_{\varepsilon_m, z_{m,i}, \hat{q}_{\varepsilon_m,i}}-\hat{q}_{\varepsilon_m,i}|\ln\varepsilon_m|\right)^{p-1}_+\frac{\partial V_{\varepsilon_m, z_{m,i}, \hat{q}_{\varepsilon_m,i}}}{\partial x_\hbar}-  p\int_{\Omega}u_m\left (V_{\varepsilon_m, Z_m}-q|\ln\varepsilon_m|\right )^{p-1}_+\left( \eta_i\frac{\partial V_{\varepsilon_m, Z_{m},i}}{\partial x_\hbar}\right)\\
=&p\int_{\Omega}u_m\left( V_{\varepsilon_m, z_{m,i}, \hat{q}_{\varepsilon_m,i}}-\hat{q}_{\varepsilon_m,i}|\ln\varepsilon_m|\right)^{p-1}_+\frac{\partial V_{\varepsilon_m, z_{m,i}, \hat{q}_{\varepsilon_m,i}}}{\partial x_\hbar}\\
&-p\int_{\Omega}u_m\left( V_{\varepsilon_m, z_{m,i}, \hat{q}_{\varepsilon_m,i}}-\hat{q}_{\varepsilon_m,i}|\ln\varepsilon_m|+O\left( \varepsilon_m^\gamma\right) \right)^{p-1}_+\frac{\partial V_{\varepsilon_m, z_{m,i}, \hat{q}_{\varepsilon_m,i}}}{\partial x_\hbar}+O\left( \varepsilon_m^{1+\gamma}\right)\\
=&O\left( \varepsilon_m^{1+\gamma}\right).
\end{split}
\end{equation}
As for the remaining terms of \eqref{301}, %from the choice of $ \eta_i $, we have 
using $ ||\nabla\eta_i||_{L^\infty}\leq C|\ln\varepsilon_m|^{2}$ and $ ||\nabla^2\eta_i||_{L^\infty}\leq C|\ln\varepsilon_m|^{4} $ we have
\begin{equation}\label{304}
\begin{array}{ll}
\int_{\Omega}\eta_iu_m\varepsilon_m^2\text{div}\left (\frac{\partial K(x)}{\partial x_\hbar}\nabla V_{\varepsilon_m, Z_{m},i}  \right )&\\
\qquad=\int\limits_{|T_{z_{m,i}}(x-z_{m,i})|\leq s_{\varepsilon_m,i}}\eta_iu_m \left( \varepsilon_m^2\text{div}\left( \frac{\partial K(x)}{\partial x_\hbar}\nabla V_{\varepsilon_m, z_{m,i}, \hat{q}_{\varepsilon_m,i}}\right)\right)&\\
\qquad\,\,\,\,\,+\int\limits_{ s_{\varepsilon_m,i}<|T_{z_{m,i}}(x-z_{m,i})|\leq |\ln\varepsilon_m|^{2}}
\eta_iu_m \left( \varepsilon_m^2\text{div}\left( \frac{\partial K(x)}{\partial x_\hbar}\nabla V_{\varepsilon_m, z_{m,i}, \hat{q}_{\varepsilon_m,i}}\right)\right)+O\left( \varepsilon_m^2\right)  &\\
\qquad=O\left(\varepsilon_m^2\right) +O\left (\varepsilon_m^2|\ln\varepsilon_m|\right )+O\left( \varepsilon_m^2\right)&\\
\qquad=O\left( \varepsilon_m^2|\ln\varepsilon_m|\right),
\end{array}
\end{equation}
\begin{equation}\label{305}
\begin{split}
-2\int_{\Omega}&u_m\varepsilon_m^2\left (K(x)\nabla\eta_i|\nabla   \frac{\partial V_{\varepsilon_m, Z_{m},i}}{\partial x_\hbar}  \right )\\
=&-2\int_{B_{|\ln\varepsilon_m|^{2}}(z_{m,i})\backslash B_{\frac{|\ln\varepsilon_m|^{2}}{2}}(z_{m,i})}u_m\varepsilon_m^2\left (K(x)\nabla\eta_i|\nabla   \frac{\partial V_{\varepsilon_m, z_{m,i},\hat{q}_{\varepsilon_m,i}}}{\partial x_\hbar}  \right )+O\left( \varepsilon_m^2|\ln\varepsilon_m|^{2}\right)\\
=&O\left( \varepsilon_m^2|\ln\varepsilon_m|^{2}\right),
\end{split}
\end{equation}
and
\begin{equation}\label{306}
\begin{split}
-\int_{\Omega}&u_m\varepsilon_m^2\text{div}\left (K(x)\nabla   \eta_i  \right )\frac{\partial V_{\varepsilon_m, Z_{m},i}}{\partial x_\hbar}\\
=&-\int_{B_{|\ln\varepsilon_m|^{2}}(z_{m,i})\backslash B_{\frac{|\ln\varepsilon_m|^{2}}{2}}(z_{m,i})}u_m\varepsilon_m^2\text{div}\left (K(x)\nabla   \eta_i  \right )\frac{\partial V_{\varepsilon_m, z_{m,i},\hat{q}_{\varepsilon_m,i}}}{\partial x_\hbar}+O\left( \varepsilon_m^2|\ln\varepsilon_m|^{4}\right)\\
=&O\left( \varepsilon_m^2|\ln\varepsilon_m|^{4}\right),
\end{split}
\end{equation}
where we have used \eqref{200}. Inserting \eqref{302}, \eqref{304}, \eqref{305} and \eqref{306} into \eqref{301}, we finally get
\begin{equation*}
\begin{split}
&\int_{\Omega}u_mL_{\varepsilon_m}\left( \eta_i\frac{\partial V_{\varepsilon_m, Z_{m},i}}{\partial x_\hbar}\right) =O\left( \varepsilon_m^{1+\gamma}\right),
\end{split}
\end{equation*}
which combined with \eqref{coef of C} deduces
\begin{equation*}
b_{j,\tilde{h},m}=O\left (\varepsilon_m^{1+\gamma}\right ).
\end{equation*}
Consequently
\begin{equation*}
\begin{split}
\sum_{j=1}^N\sum_{\tilde{h}=1}^2&b_{j,\tilde{h},m}\left( -\varepsilon_m^2\text{div}\left (K(z_{m,j})\nabla   \frac{\partial V_{\varepsilon_m, z_{m,j},\hat{q}_{\varepsilon_m,j}}}{\partial x_{\tilde{h}}}\right )\right)\\
%=&p\sum_{j=1}^m\sum_{h=1}^2b_{j,h,N}(V_{\delta_N, z_{N,j}, \hat{q}_{\delta_N,j}}-\hat{q}_{\delta_N,j})^{p-1}_+ \frac{\partial V_{\delta_N, z_{N,j},\hat{q}_{\delta_N,j}}}{\partial x_h} \\
&=O\left( \sum_{j=1}^N\sum_{\tilde{h}=1}^2\varepsilon_m^{\frac{2}{p}-1}|b_{j,\tilde{h},m}|\right)
=O\left( \varepsilon_m^{\frac{2}{p}+\gamma}\right),\ \ \ \ \text{in}\ L^p(\Omega).
\end{split}
\end{equation*}
Taking this into \eqref{208}, we have
\begin{equation}\label{307}
\begin{split}
L_{\varepsilon_m} u_m=&Q_{\varepsilon_m} L_{\varepsilon_m} u_m+\sum_{j=1}^N\sum_{\tilde{h}=1}^2b_{j,\tilde{h},m}\left( -\varepsilon_m^2\text{div}\left (K(z_{m,j})\nabla   \frac{\partial V_{\varepsilon_m, z_{m,j},\hat{q}_{\varepsilon_m,j}}}{\partial x_{\tilde{h}}}\right )\right)\\
=&O\left( \frac{1}{m}\varepsilon_m^{\frac{2}{p}}\right) +O\left(  \varepsilon_m^{\frac{2}{p}+\gamma}\right)
=o\left( \varepsilon_m^{\frac{2}{p}}\right),\ \ \ \ \text{in}\ L^p(\Omega).
\end{split}
\end{equation}

For fixed $ i, $ we define the scaled function $ \tilde{u}_{m,i}(y)=u_m(s_{\varepsilon_m, i}y+z_{m,i}) $ for $ y\in \Omega_{m,i}:=\{y\in\mathbb{R}^2\mid s_{\varepsilon_m, i}y+z_{m,i}\in \Omega\} $.
Define
\begin{equation*}
\begin{split}
\tilde{L}_{m,i}u=&-\text{div}\left (K(s_{\varepsilon_m, i}y+z_{m,i})\nabla u\right )\\
&-p\frac{s_{\varepsilon_m,i}^2}{\varepsilon_m^2}\left (V_{\varepsilon_m, Z_m}(s_{\varepsilon_m, i}y+z_{m,i})-q(s_{\varepsilon_m, i}y+z_{m,i})|\ln\varepsilon_m|\right )^{p-1}_+ u.
\end{split}
\end{equation*}
Then
\begin{equation*}
\begin{split}
||\tilde{L}_{m,i}\tilde{u}_{m,i}||_{L^p(\Omega_{m,i})}
%=&\left( \int_{\Omega_{N,i}}\left( \tilde{L}_{N,i}\tilde{u}_{N,i} \right)^pdy\right)^{\frac{1}{p}}  \\
%=&\left( \frac{1}{s_{\delta_N, i}^2}\int_{\Omega}\left( -s_{\delta_N, i}^2\text{div}(K(x)\nabla u_N)-p\frac{s_{\delta_N, i}^2}{\delta_N^2}(V_{\delta_N,Z_N}-q)^{p-1}_+u_N \right)^pdx\right)^{\frac{1}{p}}\\
=\frac{s_{\varepsilon_m, i}^2}{s_{\varepsilon_m, i}^{\frac{2}{p}}\varepsilon_m^2}||L_{\varepsilon_m}u_m||_{L^p(\Omega)}.
\end{split}
\end{equation*}
%i.e., $ s_{\delta_N, i}^{\frac{2}{p}}\cdot\frac{\delta_N^2}{s_{\delta_N, i}^2}||\tilde{L}_{N,i}\tilde{u}_{N,i}||_{L^p(\Omega_{N,i})}=||L_{\delta_N}u_N||_{L^p(\Omega)}. $
Since  $ s_{\varepsilon_m, i}=O(\varepsilon_m) $, by \eqref{307} we have
\begin{equation*}
\tilde{L}_{m,i}\tilde{u}_{m,i}=o(1)\ \ \ \  \text{in}\ \ L^p(\Omega_{m,i}).
\end{equation*}
Since $ ||\tilde{u}_{m,i}||_{L^\infty(\Omega_{m,i})}=1 $, by the classical regularity theory of  elliptic equations, $ \tilde{u}_{m,i} $ is uniformly bounded in $ W^{2,p}_{loc}(\mathbb{R}^2) $, which implies that there is $u_i$ such that
\begin{equation*}
\tilde{u}_{m,i}\to u_i\ \ \ \ \text{in}\ \ C^1_{loc}(\mathbb{R}^2).
\end{equation*}

We claim that $ u_i\equiv 0. $ On the one hand, in our ansatz for $ Z\in \Lambda_{\varepsilon, N} $, $ |z_i-z_j|\geq |\ln\varepsilon|^{-1} $. So by \eqref{203}, $ z_{m,i}\to z_i $ as $ m\to\infty $ and the fact that $ \lim_{\varepsilon\to 0}\varepsilon|\ln\varepsilon|=0 $, we get
\begin{equation*}
\begin{split}
\frac{s_{\varepsilon_m,i}^2}{\varepsilon_m^2}&\left (V_{\varepsilon_m, Z_m}(s_{\varepsilon_m, i}y+z_{m,i})-q\left (s_{\varepsilon_m, i}y+z_{m,i}\right )|\ln\varepsilon_m|\right )^{p-1}_+\\
=&\frac{s_{\varepsilon_m,i}^2}{\varepsilon_m^2}\left( V_{\varepsilon_m, z_{m,i}, \hat{q}_{\varepsilon_m,i}}(s_{\varepsilon_m, i}y+z_{m,i})-\hat{q}_{\varepsilon_m,i}|\ln\varepsilon_m|+O\left (\varepsilon_m^{\gamma}\right )\right)^{p-1}_+ \\
\to&\phi(T_{z_i}y)^{p-1}_+\ \  \  \text{in}\ C^0_{loc}(\mathbb{R}^2)\ \ \ \ \ \ \  \text{as}\ m\to\infty.
\end{split}
\end{equation*}
So $ u_i $ satisfies
\begin{equation*}
-\text{div}(K(z_i)\nabla u_i(x))-p\phi(T_{z_i}x)^{p-1}_+u_i(x)=0,\ \ x\in\mathbb{R}^2.
\end{equation*}
Let $ \hat{u}_i(x)=u_i(T_{z_i}^{-1}x) $. Since $ T_{z_i}^{-1}(T_{z_i}^{-1})^T=K(z_i) $, we have
\begin{equation*}
-\Delta \hat{u}_i(x)=-\text{div}(K(z_i)\nabla u_i)(T_{z_i}^{-1}x)=p\phi(x)^{p-1}_+\hat{u}_i(x),\ \ \ \   \ x\in \mathbb{R}^2.
\end{equation*}
By Proposition \ref{Non-degenerate}, there exist $ c_1,c_2 $ such that
\begin{equation}\label{210}
\hat{u}_i=c_1\frac{\partial \phi}{\partial x_1}+c_2\frac{\partial \phi}{\partial x_2}.
\end{equation}

On the other hand, since $ u_m\in E_{\varepsilon_m,Z_m} $, we get
\begin{equation*}
\int_{\Omega}-\varepsilon_m^2\text{div}\left( K(x)\nabla \left( \eta_i\frac{\partial V_{\varepsilon_m, Z_{m}, i}}{\partial x_\hbar}\right)\right)  u_m=0,\ \ \forall\ i=1,\cdots,N,\  \hbar=1,2,
\end{equation*}
which implies that
\begin{equation}\label{209}
\begin{split}
0%=&-\int_{\Omega}u_N\eta_i\delta_N^2\text{div}\left( K(x)\nabla \frac{\partial V_{\delta_N, z_{N,i}, \hat{q}_{\delta_N,i}}}{\partial x_\hbar}\right) -2\int_{\Omega}u_N\delta_N^2\left (K(x)\nabla\eta_i|\nabla   \frac{\partial V_{\delta_N, z_{N,i},\hat{q}_{\delta_N,i}}}{\partial x_\hbar}  \right )\\
%& -\int_{\Omega}u_N\delta_N^2\text{div}\left (K(x)\nabla   \eta_i  \right )\frac{\partial V_{\delta_N, z_{N,i},\hat{q}_{\delta_N,i}}}{\partial x_\hbar}\\
=&p\int_{\Omega}u_m\left( V_{\varepsilon_m, z_{m,i}, \hat{q}_{\varepsilon_m,i}}-\hat{q}_{\varepsilon_m,i}|\ln\varepsilon_m|\right)^{p-1}_+\frac{\partial V_{\varepsilon_m, z_{m,i}, \hat{q}_{\varepsilon_m,i}}}{\partial x_\hbar}+\int_{\Omega}\eta_iu_m\varepsilon_m^2\text{div}\left (\frac{\partial K(x)}{\partial x_\hbar}\nabla V_{\varepsilon_m, Z_{m},i}  \right )\\
&-2\int_{\Omega}u_m\varepsilon_m^2\left (K(x)\nabla\eta_i|\nabla   \frac{\partial V_{\varepsilon_m, Z_{m},i}}{\partial x_\hbar}  \right ) -\int_{\Omega}u_m\varepsilon_m^2\text{div}\left (K(x)\nabla   \eta_i  \right )\frac{\partial V_{\varepsilon_m, Z_{m},i}}{\partial x_\hbar}.
\end{split}
\end{equation}
Using again \eqref{304}, \eqref{305} and \eqref{306}, we have
\begin{equation}\label{209-1}
\begin{split}
\int_{\Omega}&\eta_iu_m\varepsilon_m^2\text{div}\left (\frac{\partial K(x)}{\partial x_\hbar}\nabla V_{\varepsilon_m, Z_{m},i}  \right )-2\int_{\Omega}u_m\varepsilon_m^2\left (K(x)\nabla\eta_i|\nabla   \frac{\partial V_{\varepsilon_m, Z_{m},i}}{\partial x_\hbar}  \right ) \\
&-\int_{\Omega}u_m\varepsilon_m^2\text{div}\left (K(x)\nabla   \eta_i  \right )\frac{\partial V_{\varepsilon_m, Z_{m},i}}{\partial x_\hbar}
=O\left( \varepsilon_m^2|\ln\varepsilon_m|^{4}\right).
\end{split}
\end{equation}
%and
%\begin{equation}\label{209-2}
%-2\int_{\Omega}u_N\delta_N^2\left (K(x)\nabla\eta_i|\nabla   \frac{\partial V_{\delta_N, z_{N,i},\hat{q}_{\delta_N,i}}}{\partial x_\hbar}  \right ) -\int_{\Omega}u_N\delta_N^2\text{div}\left (K(x)\nabla   \eta_i  \right )\frac{\partial V_{\delta_N, z_{N,i},\hat{q}_{\delta_N,i}}}{\partial x_\hbar}=O\left( \frac{\varepsilon_N^2}{|\ln\varepsilon_N|^{p}}\right).
%\end{equation}
From  \eqref{200},
\begin{equation}\label{209-3}
\begin{split}
&p\int_{\Omega} u_m\left (V_{\varepsilon_m, z_{m,i}, \hat{q}_{\varepsilon_m,i}}-\hat{q}_{\varepsilon_m,i}|\ln\varepsilon_m|\right )^{p-1}_+ \frac{\partial V_{\varepsilon_m, z_{m,i}, \hat{q}_{\varepsilon_m,i}}}{\partial x_\hbar}\\
=&p\int_{\Omega}\frac{1}{s_{\varepsilon_m,i}}\left( \frac{\varepsilon_m}{s_{\varepsilon_m,i}}\right)^{\frac{2p}{p-1}} \phi\left( \frac{T_{z_{m,i}}(x-z_{m,i})}{s_{\varepsilon_m,i}}\right)^{p-1}_+\phi'\left( \frac{T_{z_{m,i}}(x-z_{m,i})}{s_{\varepsilon_m,i}}\right)\frac{(T_{z_{m,i}})_\hbar^T\cdot T_{z_{m,i}}(x-z_{m,i})}{|T_{z_{m,i}}(x-z_{m,i})|}u_m\\
=&ps_{\varepsilon_m,i}\left( \frac{\varepsilon_m}{s_{\varepsilon_m,i}}\right)^{\frac{2p}{p-1}}\int_{\mathbb{R}^2}\phi(T_{z_{m,i}}y)^{p-1}_+\phi'(T_{z_{m,i}}y)\frac{(T_{z_{m,i}})_\hbar^T\cdot T_{z_{m,i}}y}{|T_{z_{m,i}}y|}\tilde{u}_{m,i}(y)dy.
\end{split}
\end{equation}
Taking \eqref{209-1}  and \eqref{209-3} into \eqref{209}
%, we have
%\begin{equation}\label{209-0}
%\begin{split}
%0=&ps_{\delta_N,i}\left( \frac{\delta_N}{s_{\delta_N,i}}\right)^{\frac{2p}{p-1}}\int_{\mathbb{R}^2}\phi(T_{z_{N,i}}y)^{p-1}_+\phi'(T_{z_{N,i}}y)\frac{(T_{z_{N,i}})_\hbar^t\cdot T_{z_{N,i}}y}{|T_{z_{N,i}}y|}\tilde{u}_{N,i}(y)dy+O\left( \delta_N^2|\ln\varepsilon_N|^{2M+1}\right).
%\end{split}
%\end{equation}
, dividing both sides of \eqref{209} into $ ps_{\varepsilon_m,i}(\frac{\varepsilon_m}{s_{\varepsilon_m,i}})^{\frac{2p}{p-1}} $ and passing $ m $ to the limit, we get
\begin{equation*}
\begin{split}
0%=&\int_{\mathbb{R}^2}\phi(T_{z_{i}}y)^{p-1}_+\phi'(T_{z_{i}}y)\frac{(T_{z_{i}})_\hbar^t\cdot T_{z_{i}}y}{|T_{z_{i}}y|}u_{i}(y)dy\\
=\int_{\mathbb{R}^2}\phi(x)^{p-1}_+\phi'(x)\frac{(T_{z_{i}})_\hbar^T\cdot x}{|x|}\hat{u}_{i}(x)\sqrt{\det(K(z_i))}dx,\ \  \hbar=1,2,
\end{split}
\end{equation*}
which implies that
\begin{equation}\label{211}
0=\int_{B_1(0)}\phi^{p-1}_+\frac{\partial \phi}{\partial x_\hbar}\hat{u}_i.
\end{equation}
So $ c_1=c_2=0. $ That is, $ u_i\equiv 0. $ We conclude that $ \tilde{u}_{N,i}\to 0 $ in  $ C^1(B_L(0)) $, which implies that
\begin{equation}\label{212}
||u_m||_{L^\infty(B_{Ls_{\varepsilon_m,i}}(z_{m,i}))}=o(1).
\end{equation}

Since $ Q_{\varepsilon_m} L_{\varepsilon_m} u_m = 0 $ in $ \Omega \backslash\cup_{i=1}^NB_{Ls_{\varepsilon_m,i}}(z_{m,i}) $, by the definition of $  Q_{\varepsilon_m} $ we have for $ L $ large
\begin{equation*}
 L_{\varepsilon_m} u_m = 0\  \ \text{in} \ \Omega \backslash\cup_{i=1}^NB_{Ls_{\varepsilon_m,i}}(z_{m,i}).
\end{equation*}
From Lemma \ref{lemA-5},  $ (V_{\varepsilon_m,Z_m}-q|\ln\varepsilon_m|)_+=0\  \ \text{in} \ \Omega \backslash\cup_{i=1}^NB_{Ls_{\varepsilon_m,i}}(z_{m,i}). $
So $ -\text{div}(K(x)\nabla u_m)=0 $ in $ \Omega \backslash\cup_{i=1}^NB_{Ls_{\varepsilon_m,i}}(z_{m,i}). $
%Since $ u_N=o(1) $ on $ \cup_{i=1}^m\partial B_{Ls_{\delta_N,i}}(z_{N,i}) $ and $ u_N=0 $ on $ \partial \Omega $,
Thus by the maximum principle, 
\begin{equation*}
||u_m||_{L^\infty(\Omega\backslash \cup_{i=1}^NB_{Ls_{\varepsilon_m,i}}(z_{m,i}))}=o(1),
\end{equation*}
which combined with \eqref{212} leads to
\begin{equation*}
||u_m||_{L^\infty(\Omega)}=o(1).
\end{equation*}
This is a contradiction since $ ||u_m||_{L^\infty(\Omega)}=1. $



\end{proof}

\end{document}
